\section{Einleitung}
Diese Arbeit beschreibt eine reduzierte Predictionsmarkt-DApp. Die Anwendung
verbindet einen ERC-20-Demonstrationstoken, marktbasierte Einsaetze auf diskrete
Ausgaenge und eine administrative Aufloesung nach Fristende. Im Mittelpunkt steht
damit kein allgemeines Blockchain-Geruest, sondern ein klar abgegrenzter
Kernfluss mit Tokenverteilung, Markterstellung, Einsatz und Auszahlung.

Als Beispiel dient ein Predictionsmarkt. Das Beispiel stellt fachliche und
technische Fragen kompakt dar. Die Teilnahme erfordert keine tiefe Vorarbeit.
Teilnehmende koennen direkt an einem binären Markt mitwirken. Gleichzeitig sind
Ergebnisfindung, Vertrauensannahmen und Informationsasymmetrien zentrale
Aspekte.

Das übergeordnete Ziel besteht darin, eine durchgängig lauffaehige DApp
bereitzustellen, deren Aufbau, Vertragslogik und Kopplung zur Oberfläche
nachvollziehbar bleiben. Konkret adressiert die Implementierung folgende Leitfragen:
\begin{itemize}
  \item Wie lassen sich Markterstellung, Teilnahme und Abrechnung so in Smart
    Contracts formulieren, dass der Zustand vollständig on-chain nachvollziehbar
    bleibt?
  \item Wie gestaltet sich die Kopplung zwischen Browser-Oberflaeche, Wallet und
    JSON-RPC-Endpunkt, sodass typische Nutzeraktionen (Genehmigen, Staken,
    Auflösen, Beanspruchen) robust abgebildet werden?
  \item Welche Sicherheits- und Vertrauensannahmen sind für eine Demo akzeptabel,
    und wo müsste ein produktives System nachrüsten?
\end{itemize}

Ziel dieses Berichts ist eine kurze, fachlich geschlossene technische
Dokumentation im Stil einer Projektarbeit. Sie beschreibt die Lösung, analysiert
die wichtigsten Designentscheidungen und leitet daraus die Grenzen und das
Verbesserungspotenzial der Implementierung ab.

Kapitel~\ref{chap:theorie} beschreibt die Methodik und die technische Grundidee.
Kapitel~\ref{chap:kern} analysiert den Kernablauf der DApp und den
Funktionsumfang. Kapitel~\ref{chap:sicherheit} diskutiert Grenzen und
Verbesserungspotenzial. Zusammen ergibt sich daraus ein kurzer, geschlossen
argumentierter Projektbericht.

Die Arbeit ist als Demonstrationssystem konzipiert und nicht für den Einsatz mit
realen Daten vorgesehen. Die folgenden Punkte sind daher keine nachträglichen
Defizite, sondern bewusst gesetzte Grenzen des Projektumfangs.

Als Sicherheiten-Token dient ausschliesslich der eigene ERC-20-Token \emph{VOTE} mit
einmaligem Faucet. Jede Adresse kann sich einmalig einen Betrag zuweisen. Damit
werden weder natives Ether noch ein Stablecoin als Einsatzmittel verwendet.
Preisrisiko, Compliance und Brueckenlogik bleiben damit ausserhalb des Systems.
Brueckenlogik bezeichnet hier die Komplexitaet, die sich bei Abhaengigkeiten zu
anderen Systemen oder zum Mainnet ergibt.

Die Aufloesung erfolgt durch eine beim Deployment gesetzte \emph{Admin}-Adresse.
Diese Rolle bestimmt nach Fristende den gewinnenden Ausgang. Das System enthält
keine Chainlink-Feeds, keine Reporter-Jury und keine Einspruchsinstanz. Chainlink-
Feeds sind externe Datenquellen, die in einem realen Anwendungsfall zur Aufloesung
eines Markts noetig sein koennen.

der Vertrag \texttt{PredictionMarket} erlaubt formal mehrere Ausgaenge, sofern
die Validierung in \texttt{createMarket} erfuellt ist. Die Nuxt-Oberfläche
vereinfacht die Nutzerfuehrung, indem sie neue Maerkte derzeit auf die beiden
festen Ausgaenge \emph{YES} und \emph{NO} festlegt. Die Benutzeroberfläche bildet
damit bewusst einen binären Kernfall ab.

Im weiteren Verlauf bezeichnet der Begriff \emph{Predictionsmarkt} die konkret
implementierte Variante: Ein Markt besitzt eine textuelle Frage, einen
Aufloesungskontext zur operationalen Klaerung, eine absolute Deadline, eine
Menge benannter Ausgaenge und pro Ausgang einen Token-Pool. Solange die Deadline
nicht überschritten ist, koennen Adressen Token in die Pools einbringen
(\emph{staken}). Nach Fristende wählt die Admin-Adresse den gewinnenden Ausgang.
Gewinner beanspruchen mit einer eigenstaendigen Transaktion einen anteiligen
Anteil am gesamten Markt-Pool gemaess der im Vertrag implementierten Formel.

Die aus den Poolverhaeltnissen abgeleiteten Prozentwerte sind in diesem Kontext
\textbf{einsatzgewichtete Anteile}. Sie dienen dem Nutzenden als Feedback fuer
einen freiwillig platzierten Einsatz. Es handelt sich somit nicht um eine exakte
Vorhersage eines realen Ereignisses.
